\documentclass[11pt,a4paper]{article}
\usepackage[margin=25mm]{geometry}
\usepackage[T1]{fontenc}
\usepackage{lmodern}
\usepackage{amsmath}
\usepackage{booktabs}
\usepackage{graphicx}
\usepackage{hyperref}
\hypersetup{colorlinks=true,linkcolor=blue,urlcolor=blue}
\setlength{\parskip}{0.5em}
\setlength{\parindent}{0pt}

\title{Korean Inflation Forecasts\\Results and Method}
\author{Model output review}
\date{16 September 2026\\Input snapshot through June 2026}

\begin{document}
\maketitle

\section{Results}
The saved monthly vector autoregressions (VARs) project headline CPI, core CPI excluding food and energy, and producer prices (PPI) for 36 months. Headline year-ended CPI inflation declines from 3.17\% in June 2026 to 2.23\% one year later, then reaches 2.56\% in June 2029. Core CPI rises to 2.91\% in June 2027 before easing to 2.37\% by the end of the path. PPI inflation falls sharply from 8.42\% in the last observed month to 2.13\% in June 2027 and 1.90\% in June 2029. The PPI decline is a model projection, not an assumption imposed on the CPI systems.

\begin{center}
\begin{tabular}{lrrr}
\toprule
Period & Headline CPI & Core CPI & PPI \\
\midrule
June 2026 (observed) & 3.17 & 2.48 & 8.42 \\
June 2027 & 2.23 & 2.91 & 2.13 \\
June 2028 & 2.63 & 2.51 & 1.87 \\
June 2029 & 2.56 & 2.37 & 1.90 \\
\bottomrule
\end{tabular}
\end{center}
\noindent All entries are year-ended percentage changes, rounded from the Korean CPI and PPI forecast CSVs. Solid chart lines are observed rates; dashed lines are projections. A dotted rule marks June 2026.

\begin{figure}[h]
\centering
\includegraphics[width=\linewidth]{\detokenize{../results/figures/kr_cpi_inflation_forecast.png}}
\caption{Korean headline and core CPI inflation and 36-month forecasts.}
\end{figure}

\clearpage
\begin{figure}[h]
\centering
\includegraphics[width=\linewidth]{\detokenize{../results/figures/kr_ppi_inflation_forecast.png}}
\caption{Korean producer-price inflation and 36-month forecast.}
\end{figure}

\clearpage
\section{Methodology}
\subsection{Data and specification}
The country workbook contains only the three target price indices and four drivers required from the original \texttt{KR} sheet. The targets are headline CPI, core CPI excluding food and energy, and PPI. Each has a separate monthly VAR with the same driver block: inflation expectations, WTI oil, import prices, and unemployment. Target, oil, and import-price indices enter as 100 times their month-on-month percentage changes; expectations and unemployment enter as levels. The code begins assembling data in January 2002; after alignment and differencing, the transformed estimation samples start in March 2002 and end in June 2026.

Each fitted system has the form
\[
  x_t=c+A_1x_{t-1}+\cdots+A_px_{t-p}+u_t.
\]
\texttt{statsmodels} selects $p$ by the Akaike information criterion from at most 12 lags. The saved headline CPI and PPI fits use three lags; core CPI uses six. The VAR recursively projects the target and all four drivers for 36 months. Projected monthly target changes are chained to the last observed index; the reported year-ended rate is $100(P_t/P_{t-12}-1)$.

Run \texttt{py -3.14 run\_all\_var\_models.py} from the \texttt{macro-work} folder to regenerate the CSVs and charts from the bundled inputs. The CSVs store historical and forecast index levels and year-ended rates.

\section{Interpretation and limits}
Separate VARs mean the CPI and PPI paths are not constrained to agree with one another. In particular, PPI is a forecast target in its own model, not a forecast driver in the CPI models. The charts show point forecasts only; no uncertainty bands or rolling out-of-sample evaluation are published. The fitted models extrapolate historical relationships and do not encode future policy decisions or new commodity-price shocks. The rapid PPI disinflation should therefore be read as a feature of this statistical forecast, with uncertainty around its timing and size.

\end{document}
